\section{Operational model}
\label{sec:model}

Fix $n\geq1$. At slot $i$, a trusted verifier prepares
\[
\ket{\Phi}_{R_iA_i}=\frac{\ket{00}+\ket{11}}{\sqrt2},
\]
sends $A_i$ to the device, receives $B_i$, and sequesters $B_i$ before releasing $A_{i+1}$. After the last output, the device commits a terminal memory port $M$ to the verifier. A hidden coin then selects one branch.

\paragraph{AUDIT.}
The verifier measures every $R_i$ in $Z$, obtaining independent uniform bits $X_i$, and decodes from $M$ a guess for
\[
Z=X_1\oplus\cdots\oplus X_n.
\]
The sequestered $B_i$ are unavailable. Let $\PA$ be the success probability.

\paragraph{RETURN.}
The verifier applies a fixed joint decoder to $M B_1\cdots B_n$ and accepts with
\[
\Pi_{\mathrm R}=\proj{0}_M\otimes
\bigotimes_{i=1}^n\proj{\Phi}_{R_iB_i}.
\]
Let $\FR$ be the acceptance probability. The visible reset $M=0$ does not imply deletion of any inaccessible purification of a classical record.

For $0\leq\lambda\leq1$, define the support score
\[
\Sscore=\lambda\PA+(1-\lambda)\FR.
\]
The balanced benchmark uses $\lambda=1/2$.

\subsection{Adaptive classical-memory streaming null}

The null may use arbitrary finite classical memory and shared randomness. In slot $i$, conditional on the earlier classical history $h=k_{<i}$, it applies an arbitrary instrument
\[
\{\mathcal N^{(i,h)}_{k_i}:A_i\rightarrow B_i\}_{k_i}.
\]
It may create and discard arbitrary within-slot quantum ancillas. It may not retain a coherent system across the slot boundary, begin with entanglement shared between slot ancillas, revisit a committed $B_i$, or receive a $B_i$ in AUDIT. In RETURN, the full classical transcript may control an arbitrary joint recovery on all sequestered outputs.

After refining all classical random seeds and outcomes, every leaf transcript $k=(k_1,\ldots,k_n)$ has a conditional channel that factorises over slots. For computational-basis inputs $x=(x_1,\ldots,x_n)$, its likelihood is
\begin{equation}
p(k|x)=\prod_{i=1}^n p_i(k_i|x_i,k_{<i}).
\label{eq:factor}
\end{equation}
This factorisation is the mathematical null. The experiment must support it through timing and isolation evidence; the score alone cannot enforce a laboratory boundary.

\subsection{Local audit strength and return curve}

At a realised prefix $h$, define
\[
r_i(k_i|h)=\frac{p_i(k_i|0,h)+p_i(k_i|1,h)}{2}
\]
and, when $r_i>0$, define a signed posterior bias by
\begin{equation}
s_i\delta_i=
\frac{p_i(k_i|0,h)-p_i(k_i|1,h)}
{p_i(k_i|0,h)+p_i(k_i|1,h)},
\qquad 0\leq\delta_i\leq1.
\label{eq:bias}
\end{equation}
The one-slot return curve is
\begin{equation}
f(t)=\frac{1+\sqrt{1-t^2}}{2}.
\label{eq:fcurve}
\end{equation}
Here $t=0$ extracts no computational-basis information and preserves an EPR pair perfectly; $t=1$ permits a perfect bit record and leaves entanglement fidelity $1/2$.

\subsection{A saturating null family}

For any $t_i\in[0,1]$, the binary diagonal instrument
\[
K^{(i)}_0=
\begin{pmatrix}
\sqrt{(1+t_i)/2}&0\\0&\sqrt{(1-t_i)/2}
\end{pmatrix},\quad
K^{(i)}_1=
\begin{pmatrix}
\sqrt{(1-t_i)/2}&0\\0&\sqrt{(1+t_i)/2}
\end{pmatrix}
\]
produces local audit bias $t_i$ and return fidelity $f(t_i)$. The transcript parity is the optimal AUDIT guess and identity recovery saturates the RETURN expression derived below.
